\section{Identity extensions of invariant operators}
\label{sec:projection-identity-extension}

The transport metrics in \cite[Section~7]{OpenAI2026AreaLaw} act on the
symmetric subspace. The following identities describe an extension of an
invariant operator to the ambient finite-dimensional complex Hilbert space.
Let $P$ be an orthogonal projection and write $Q=I-P$. For an operator $A$
commuting with $P$, put
\begin{align*}
 E_P(A)&=PAP+Q.
\end{align*}
Its restriction to the complementary subspace is the identity. Lower
operator bounds are separate assertions.

\begin{lemma}[Restriction of the identity extension]
\label{lem:projection-identity-extension-restriction}
\lean{Matrix.compression_add_one_sub_mul_projection}
\leanok
If $AP=PA$, then
\begin{align*}
 E_P(A)P&=AP.
\end{align*}
\end{lemma}

\begin{proof}
\leanok
Use $P^2=P$, $QP=0$, and $PAP=AP$.
\end{proof}

\begin{lemma}[Products of identity extensions]
\label{lem:projection-identity-extension-product}
\lean{Matrix.compression_add_one_sub_mul_compression_add_one_sub}
\leanok
If $AP=PA$ and $BP=PB$, then
\begin{align*}
 E_P(A)E_P(B)&=ABP+Q.
\end{align*}
\end{lemma}

\begin{proof}
\leanok
The mixed terms vanish because $PQ=QP=0$. The other two terms are $ABP$
and $Q^2=Q$.
\end{proof}

\begin{lemma}[Commutation of identity extensions]
\label{lem:projection-identity-extension-commutation}
\lean{Matrix.commute_compression_add_one_sub_of_mul_projection_eq}
\uses{lem:projection-identity-extension-product}
\leanok
Suppose $AP=PA$, $BP=PB$, and $ABP=BAP$. Then $E_P(A)$ and $E_P(B)$ commute.
\end{lemma}

\begin{proof}
\leanok
Apply Lemma~\ref{lem:projection-identity-extension-product} in both orders.
\end{proof}

\begin{lemma}[Functional calculus on the invariant subspace]
\label{lem:projection-identity-extension-functional-calculus}
\lean{Matrix.cfc_compression_add_one_sub_mul_projection}
\uses{lem:projection-identity-extension-restriction}
\leanok
If $A$ is Hermitian and $AP=PA$, then every real-valued function $f$ satisfies
\begin{align*}
 f(E_P(A))P&=Pf(A).
\end{align*}
Here functional calculus is defined on the finite spectra of the two
Hermitian operators.
\end{lemma}

\begin{proof}
\leanok
The operator $E_P(A)$ is Hermitian. The restriction identity gives
$E_P(A)P=PA$. Functional calculus preserves this intertwining identity:
in eigenbases of the two Hermitian operators, an entry of $P$ can be nonzero
only between equal eigenvalues. Applying $f$ preserves equality of those
eigenvalues.
\end{proof}
